\section{How the research proceeded}
\label{sec:methods}
\label{sec:trajectory}

\subsection{A shared research budget}

Kinyon and Urban asked Codex to solve as many previously unsolved Notebook
problems as possible in 48 hours, using GAP and any further useful
software, while keeping plans, proofs and frequent local commits.
The supplied twenty-first edition was updated through September 2026.
The agent used GPT-6 Astra with \texttt{xhigh} reasoning, GAP 4.16.1,
exact Python and C/C++ programs, and internet access to mathematical
sources. It worked within requested local limits of generally 20 CPU
cores and 100\,GB of memory. The research phase used one language-model
agent; parallel computational workers handled searches and certificate
checks. Claude's independent review came later.

The shared budget allowed the agent to choose problems opportunistically
and to reuse constructions. Written state preserved promising leads,
proof drafts, failed approaches and unfinished checks between work
sessions. Nearly two hundred local commits record the progression.
The original instruction and the exact configuration, counts and
accounting appear in Appendices~\ref{app:protocol}--\ref{app:methods}.

\subsection{From a lead to an argument}

The first retained lead was the definability counterexample to 21.106.
While developing short symbolic arguments, the agent launched finite
screens for other questions. Literature frequently supplied a productive
starting point: Kida's order-96 configuration led to the non-monomial
semi-abelian group, and an exception in Tsang's computational coverage
led to the multiple-holomorph example. Other questions shared deeper
inputs, notably Amelio's geometric construction for 10.62 and 4.75.
The proofs identify which steps are imported and which deductions the
experiment supplied.

Validation changed the direction of the work. In the graph-symmetry
search for 21.52, a permutation normalizing an inner-automorphism action
was initially treated as evidence for a group automorphism. The agent
replaced that weaker test by constructing an actual automorphism and
checking its action. A program can compute its specified predicate
correctly while that predicate fails to express the mathematical question.
This example explains the value of reviewing the specification as well
as the calculation.

Source checks also changed the classification of results. The agent's
argument for 17.34 was ultimately covered by an existing
strong-amalgamation theorem. A small construction for 13.19 likewise
became a rediscovery after an earlier result was found. Those arguments
remain in Section~\ref{app:prior}, with their prior-work status, because
they help explain both the mathematics and the course of the experiment.

\subsection{Long computations and partial theorems}

For 20.100, a structural reduction made each fixed distinct-powers case
amenable to a finite certificate. Generating the \(n=7\) certificate
was only an intermediate step: checking it took approximately 26 hours.
Other work continued during that computation. The resulting fixed cases
and torsion-free theorem are presented in Section~\ref{partial:20.100};
the unrestricted question for arbitrary \(n\) remained open.

Almost Engel elements and \(p\)-Jordan bounds supplied further substantial
partials. Some broad finite searches ended with useful coverage but no
complete answer. This variety of outcomes is why we retain the proofs
and their boundaries, rather than report only a solved-problem total.
Appendix~\ref{app:chronology} gives the dated sequence, the larger search
counts and the distinction between generation, checking and closeout.
The research record ends with internal review; the separate model review
in Section~\ref{sec:external-review} tests the resulting mathematical work.
